\documentclass[11pt,a4paper]{article}
\usepackage[T1]{fontenc}
\usepackage[utf8]{inputenc}
\usepackage{lmodern}
\usepackage{microtype}
\usepackage[a4paper,margin=2.25cm]{geometry}
\usepackage{longtable,booktabs,array,calc}
\usepackage{fancyvrb}
\usepackage{xcolor}
\usepackage{hyperref}
\usepackage{xurl}
\usepackage{fancyhdr}
\usepackage{enumitem}
\usepackage{parskip}
\usepackage{titlesec}
\definecolor{ddtadark}{RGB}{30,38,48}
\definecolor{ddtablue}{RGB}{38,82,126}
\hypersetup{colorlinks=true,linkcolor=ddtablue,urlcolor=ddtablue,citecolor=ddtablue}
\pagestyle{fancy}
\fancyhf{}
\lhead{Documentation-Driven Threat Analysis}
\rhead{BA4 projection candidate R2}
\cfoot{\thepage}
\setlength{\headheight}{14pt}
\setlist[itemize]{leftmargin=1.5em,itemsep=0.2em,topsep=0.25em}
\setlist[enumerate]{leftmargin=1.7em,itemsep=0.2em,topsep=0.25em}
\titleformat{\section}{\Large\bfseries\color{ddtadark}}{}{0em}{}
\titleformat{\subsection}{\large\bfseries\color{ddtadark}}{}{0em}{}
\titleformat{\subsubsection}{\normalsize\bfseries\color{ddtadark}}{}{0em}{}
\providecommand{\tightlist}{\setlength{\itemsep}{0pt}\setlength{\parskip}{0pt}}
\setlength{\emergencystretch}{4em}
\hbadness=10000
\hfuzz=20pt
\sloppy
\begin{document}
\section{DDTA BA4 projection boundary, interpretation and coverage
candidate -
R2}\label{ddta-ba4-projection-boundary-interpretation-and-coverage-candidate---r2}

\textbf{Status:} PROVISIONAL CANDIDATE AFTER BA4-T2 / NOT CLOSED / BA4
OPEN

\textbf{Derived by:} BA4-T2 method-owned interpretation, coverage loss
and cross-projection consistency pressure test

\textbf{Repository baseline reviewed:}
\texttt{f90ef3a0bc0b7712cb8081165c28e8923aec9e2d}

\textbf{Refines:}
\texttt{BA4\_PROJECTION\_BOUNDARY\_TRACEABILITY\_CANDIDATE\_R1.md}; R1
remains immutable research history.

\textbf{Closed dependencies:} BA0 responsibility boundary; BA1
\texttt{BAReferent\ +\ BAProposition}; BA2 proposition semantics; BA3
provenance/derivation/lifecycle/change contract.

\subsection{1. Purpose}\label{1-purpose}

BA4-T2 asks whether the T1 projection lower bound survives when two
consumer methods impose intentionally incompatible local taxonomies on
the same accepted Base Analysis, when selection becomes aggressively
lossy, and when a projection explicitly surfaces non-current or
unresolved BA state.

The pressure target is not whether two diagrams can be rendered. The
target is whether the projection contract can keep four boundaries
honest at the same time:

\begin{enumerate}
\def\labelenumi{\arabic{enumi}.}
\tightlist
\item
  shared BA meaning remains the common methodology-neutral basis;
\item
  consumer interpretation remains consumer-owned even when it is
  reproducibly derived;
\item
  omission/completeness is stated relative to the projection contract
  rather than silently becoming project semantics; and
\item
  cross-projection comparison can be explained through BA trace without
  inventing a universal projection ontology.
\end{enumerate}

BA4 remains representation-independent. No graph, class model, storage
schema, rendering notation or executable rule language is mandated.

\subsection{2. Refined provisional lower-bound
contract}\label{2-refined-provisional-lower-bound-contract}

T2 preserves the T1 separation among projection definition,
baseline-scoped materialization and projection-local items, but refines
the coverage and method-interpretation responsibilities.

\begin{verbatim}
BAProjectionDescriptor                         [projection metadata; NOT BAE]
|- projectionKey                         1
|- projectionRevisionKey                 1     immutable
|- consumerPurpose                       1
|- selectionCoverageContract             1
|    |- eligibleBAScope                  1
|    |- coverageMode                     1
|    |    EXHAUSTIVE_FOR_DECLARED_SCOPE
|    |    | SELECTIVE
|    |- qualificationPolicy              1
|    `- omissionSemantics                1
|- mappingContract                       1
|- interpretationBoundary                1
`- interpretationRuleDescriptor          0..*
     |- ruleKey                          1
     |- inputRoleContract                1..*
     |- applicabilityContract            1
     |- outputKindContract               1
     `- normativeDefinition              1

BAProjectionMaterialization                    [derived artifact; NOT BAE]
|- projectionRef                         1     projectionKey@revision
|- sourceBABaselineKey                   1
`- item                                  0..*

BAProjectionItem                               [projection-local; NOT BAE]
|- projectionItemKey                     1
|- projectionOwnedKind                   1
|- traceBinding                           1..*
|    |- traceRoleKey                      1
|    `- baElementRef                      1     BA element @ BA baseline
|- sharedSemanticRendering               0..1
|- methodOwnedInterpretation             0..1
`- interpretationRuleRef                 0..1
     [required when methodOwnedInterpretation is meaning-bearing]
\end{verbatim}

The nested shapes above are semantic responsibilities. An implementation
may combine them physically.

\texttt{interpretationRuleDescriptor} is consumer/projection metadata,
not a BA3 derivation rule and not a BAE family. Its rule keys are scoped
to one immutable projection revision.

\subsection{3. T1 responsibilities that survive
unchanged}\label{3-t1-responsibilities-that-survive-unchanged}

The following T1 findings survive the stronger pressure without material
revision:

\begin{verbatim}
projection derived from accepted BA                   REQUIRED
projection is project authority                        REJECTED
immutable projection revision                          REQUIRED
meaning-bearing item -> BA trace                       REQUIRED
mandatory duplicate GovernedSourceRef                  REJECTED
role-bound multi-input trace                 REQUIRED WHEN ROLES DIFFER
projection-local item identity               REQUIRED WHEN MEANING-BEARING
projection item identity -> BAE identity                REJECTED
method-owned interpretation -> shared BA                REJECTED
BA change -> projection rebuild                         REQUIRED
second project lifecycle in projection                  REJECTED
\end{verbatim}

The authority chain remains:

\begin{verbatim}
governed documentation
    -> accepted Base Analysis
        -> projection materialization
            -> method-owned interpretation
\end{verbatim}

No arrow reverses merely because a projection is useful, reproducible or
analysis-relevant.

\subsection{4. Two intentionally incompatible method
projections}\label{4-two-intentionally-incompatible-method-projections}

T2 uses two bounded projection families over the same facial-access BA
basis. Their local taxonomies are intentionally not aligned.

\subsubsection{4.1 Projection F - flow/exposure
oriented}\label{41-projection-f---flowexposure-oriented}

Illustrative projection-owned kinds:

\begin{verbatim}
FlowParticipant
FlowExchange
ExposureAnnotation
\end{verbatim}

A BA transfer may project to a \texttt{FlowExchange}:

\begin{verbatim}
transfer(
  source      -> CameraSubsystem,
  destination -> RecognitionProcessor,
  content     -> RecognitionCapture)
\end{verbatim}

The method may additionally interpret:

\begin{verbatim}
consumeService(project, LocalConnectivity)
+
assignResponsibility[negative](
  project,
  underlyingTransport,
  ownership/management)
\end{verbatim}

as a method-owned annotation such as:

\begin{verbatim}
externally-provided-transport-exposure
\end{verbatim}

That label is not BA meaning.

\subsubsection{4.2 Projection A - assurance/responsibility
oriented}\label{42-projection-a---assuranceresponsibility-oriented}

Illustrative projection-owned kinds:

\begin{verbatim}
AssuranceSubject
AuthorityPlacement
ContractDependency
NormalizationExpectation
\end{verbatim}

The same BA service/responsibility basis may instead be interpreted as:

\begin{verbatim}
delegated-assurance-dependency
\end{verbatim}

The projection may emphasize responsibility placement, service contracts
and governed normalized results while omitting ordinary flow
interactions.

\texttt{externally-provided-transport-exposure} and
\texttt{delegated-assurance-dependency} are intentionally different
local taxonomies. Neither is the canonical projection of the other, and
neither becomes a BA semantic key.

\subsection{5. No universal projection ontology is
required}\label{5-no-universal-projection-ontology-is-required}

The two method projections can disagree about:

\begin{itemize}
\tightlist
\item
  category names;
\item
  item granularity;
\item
  whether one BA element becomes an item, annotation or supporting
  trace;
\item
  whether several BA elements are aggregated into one method item; and
\item
  which BA facts are outside their selection scope.
\end{itemize}

This difference is not itself inconsistency.

The common denominator is the accepted BA trace:

\begin{verbatim}
Projection F item
    -> BA trace set / role bindings

Projection A item
    -> BA trace set / role bindings
\end{verbatim}

Cross-projection comparison may compute trace overlap, disjointness or
shared BA basis. It does not require a new shared
\texttt{ProjectionConcept}, \texttt{ProjectionRelation} or
method-category equivalence identity.

A same-looking label in two projections does not establish shared
identity. Different labels over overlapping BA trace do not establish
contradiction.

\subsection{6. T1 trace plus descriptor is insufficient for non-trivial
method
interpretation}\label{6-t1-trace-plus-descriptor-is-insufficient-for-non-trivial-method-interpretation}

T1 already requires an immutable projection descriptor and role-bound BA
trace. T2 constructs a counterexample where one descriptor revision
contains several method-owned interpretation rules that can consume
overlapping BA inputs.

Suppose one projection revision permits both:

\begin{verbatim}
R-ext-service:
  serviceUse + responsibilityBoundary
      -> externally-provided-transport-exposure

R-assurance-gap:
  serviceUse + responsibilityBoundary + constraintSet
      -> assurance-dependency-review
\end{verbatim}

An item that records only:

\begin{verbatim}
traceBinding:
  serviceUse             -> <BA consumeService>
  responsibilityBoundary -> <BA negative responsibility>
\end{verbatim}

is not enough to identify which local interpretation rule produced its
method-owned meaning.

Therefore T2 requires:

\begin{verbatim}
methodOwnedInterpretation present
    -> interpretationRuleRef required
\end{verbatim}

The reference resolves within the immutable projection revision to an
inspectable consumer-owned rule descriptor.

This requirement is narrower than BA3 derivation:

\begin{verbatim}
BA3 derivation rule
  methodology-neutral
  may produce accepted shared BA meaning

BA4 interpretation rule
  consumer/method-owned
  produces only projection-local interpretation
\end{verbatim}

The registries must not be merged.

\subsection{7. Projection interpretation
reproducibility}\label{7-projection-interpretation-reproducibility}

For method-owned interpretation, reproducibility means an independent
reviewer can determine, from:

\begin{verbatim}
same accepted BA baseline/materialization
+ same role-bound trace basis
+ same immutable projection revision
+ same interpretation rule
\end{verbatim}

whether the same materially equivalent local interpretation is
justified, or whether the rule is not applicable.

It does not require:

\begin{itemize}
\tightlist
\item
  byte-identical projection serialization;
\item
  identical generated item IDs;
\item
  one software implementation;
\item
  one rendering layout; or
\item
  promotion of the interpretation into BA.
\end{itemize}

If the same rule revision and same role-bound BA basis permit materially
incompatible local meanings without explicit method uncertainty, the
projection rule is insufficient and must be refined inside the method
projection.

\subsection{8. Coverage needs an explicit
mode}\label{8-coverage-needs-an-explicit-mode}

T1 requires a \texttt{selectionCoverageContract}, but aggressive
lossiness exposes a missing distinction.

Two projection contracts may both declare an eligible scope such as
\texttt{current\ accepted\ transfer\ propositions}, yet mean different
things:

\begin{verbatim}
Projection X:
  show every eligible transfer

Projection Y:
  show a selected subset useful for one question
\end{verbatim}

The resulting omissions have different projection semantics.

T2 therefore closes a provisional coverage-mode distinction:

\begin{verbatim}
EXHAUSTIVE_FOR_DECLARED_SCOPE
  every BA element eligible under the declared scope/qualification
  is expected to be represented according to the mapping contract

SELECTIVE
  the projection intentionally permits omission of otherwise eligible BA
  elements; omission carries no completeness claim
\end{verbatim}

This is completeness relative to BA and the projection descriptor, not
completeness of project documentation.

\subsection{9. Exhaustive omission is a projection defect, not project
negation}\label{9-exhaustive-omission-is-a-projection-defect-not-project-negation}

For \texttt{EXHAUSTIVE\_FOR\_DECLARED\_SCOPE}:

\begin{verbatim}
eligible accepted BA element omitted
    -> projection materialization defect / incomplete build
\end{verbatim}

It does \textbf{not} mean:

\begin{verbatim}
project says element is false
project documentation is incomplete
BA element does not exist
\end{verbatim}

For \texttt{SELECTIVE}:

\begin{verbatim}
item omitted
    -> no conclusion beyond this projection's selection choice
\end{verbatim}

This distinction prevents two opposite errors:

\begin{enumerate}
\def\labelenumi{\arabic{enumi}.}
\tightlist
\item
  treating lossy omission as project absence; and
\item
  allowing a projection that claims complete BA coverage to silently
  omit eligible meaning.
\end{enumerate}

No universal query language is required. The descriptor only needs an
inspectable scope and enough construction semantics to review whether
the declared coverage mode was met.

\subsection{10. Qualification policy for diagnostic, stale and pending
BA
state}\label{10-qualification-policy-for-diagnostic-stale-and-pending-ba-state}

T2 explicitly pressure-tests a review-oriented projection that may
select:

\begin{verbatim}
DIAGNOSTIC_UNRESOLVED
STALE
PENDING_REVIEW
\end{verbatim}

The T1 rule that admissible review/freshness states belong to the
coverage contract is retained and made explicit as
\texttt{qualificationPolicy}.

A normal current-state method projection may declare, for example:

\begin{verbatim}
qualificationPolicy:
  ACCEPTED + CURRENT shared BA only
\end{verbatim}

A review projection may instead include stale or diagnostic material.

The critical invariant is:

\begin{quote}
A projection may expose a non-current or unresolved BA element, but it
must not present it as a current accepted project fact.
\end{quote}

The projection may resolve BA3 state through its \texttt{BAElementRef};
it does not need to duplicate BA3 lifecycle metadata as a second
authority.

\subsection{11. Qualification propagation to method
interpretation}\label{11-qualification-propagation-to-method-interpretation}

A method-owned interpretation may be useful even when its BA basis is
stale or diagnostic, for example to show what analysis output is
potentially affected.

However:

\begin{verbatim}
stale/unresolved BA basis
    -> method interpretation may be shown as review-oriented/local
    -> MUST NOT become current shared project truth
\end{verbatim}

T2 does not introduce a new projection lifecycle enum. The
descriptor\textquotesingle s qualification policy and interpretation
boundary must preserve the source BA qualification in the
rendered/consumed semantics.

A method tool may use UI badges such as \texttt{review},
\texttt{stale\ basis} or \texttt{diagnostic\ basis}; these are
presentation choices, not new BA or BA4 semantic state families.

\subsection{12. Cross-projection consistency
contract}\label{12-cross-projection-consistency-contract}

Cross-projection consistency is intentionally weaker than taxonomy
agreement.

For two materializations over the same BA baseline, the minimum
consistency conditions are:

\begin{enumerate}
\def\labelenumi{\arabic{enumi}.}
\tightlist
\item
  each projection identifies the same source BA baseline it actually
  consumed;
\item
  each shared semantic rendering is supportable by its own traced BA
  basis;
\item
  each method-owned interpretation is explicitly local and references
  its applicable local rule when required;
\item
  each projection satisfies its own declared coverage mode and
  qualification policy; and
\item
  comparison across projections uses BA trace/BA3 continuity rather than
  assumed method-category equivalence.
\end{enumerate}

Therefore the following can both be valid:

\begin{verbatim}
Projection F:
  externally-provided-transport-exposure

Projection A:
  delegated-assurance-dependency
\end{verbatim}

when both trace honestly to the same service/responsibility BA basis
under different method rules.

\subsection{13. What counts as a real cross-projection
contradiction}\label{13-what-counts-as-a-real-cross-projection-contradiction}

A contradiction is not merely different method interpretation.

A material projection inconsistency exists when, for example:

\begin{itemize}
\tightlist
\item
  a projection presents a shared semantic rendering that inverts BA
  polarity;
\item
  a projection claims \texttt{EXHAUSTIVE\_FOR\_DECLARED\_SCOPE} while
  omitting an eligible BA item;
\item
  an item claims current project meaning from stale/unresolved basis
  contrary to its qualification policy;
\item
  a method-owned interpretation is presented as BA/project truth; or
\item
  a projection claims to use baseline B1 while tracing to an
  unacknowledged B0 BA materialization.
\end{itemize}

These are projection-contract violations. They do not require a new BA
proposition that says two projections conflict.

\subsection{14. M1-M4 rebuild pressure across incompatible
projections}\label{14-m1-m4-rebuild-pressure-across-incompatible-projections}

\subsubsection{14.1 M1 - Ethernet to
Wi-Fi}\label{141-m1---ethernet-to-wi-fi}

BA retains the transfer and service-use meaning while replacing the
realization proposition.

Projection F may remain materially unchanged if its coverage ignores
realization technology.

Projection A may replace a local realization/assurance item because it
includes technology realization in its declared scope.

Different projection deltas are expected and consistent because both
rebuild from BA@B1 under different contracts.

\subsubsection{14.2 M2 - external transport to project-owned
transport}\label{142-m2---external-transport-to-project-owned-transport}

The BA responsibility proposition changes from negative to affirmative
and external service consumption may retire.

Projection F\textquotesingle s
\texttt{externally-provided-transport-exposure} interpretation ceases to
be applicable and disappears on rebuild.

Projection A may instead produce a project-owned authority-placement
item.

No projection item mutates the BA proposition and no cross-method common
category is required.

\subsubsection{14.3 M3 - remote to local
recognition}\label{143-m3---remote-to-local-recognition}

When the remote transfer is retired, Projection F must remove the
corresponding flow exchange if its coverage is exhaustive for current
transfers.

Projection A may retain unrelated responsibility/contract items if their
BA basis remains accepted/current.

A projection that retains the old flow only because its B0 local item
existed violates rebuild semantics.

\subsubsection{14.4 M4 - representation conflict and stale
transfer}\label{144-m4---representation-conflict-and-stale-transfer}

When BA marks the old capture-transfer meaning stale and localizes a
diagnostic, a normal current-state flow projection excludes it under its
qualification policy.

A review projection may expose:

\begin{verbatim}
stale transfer basis
+ diagnostic/unresolved representation conflict
\end{verbatim}

but must visibly preserve those qualifications.

Neither method projection may silently invent
\texttt{OpaqueEvidenceReference} project truth before governance
resolves the source inconsistency.

\subsection{15. Order/WMS responsibility-boundary
pressure}\label{15-orderwms-responsibility-boundary-pressure}

When project-owned inventory moves toward an external WMS:

\begin{itemize}
\tightlist
\item
  a flow/integration projection may emphasize new
  \texttt{consumeService} interactions and external exchanges;
\item
  an assurance/responsibility projection may emphasize moved authority,
  contract dependency and responsibility boundaries; and
\item
  internal mutation propositions retired by BA disappear only from
  projections whose selection includes them.
\end{itemize}

The projection families can therefore change in very different shapes
while remaining consistent with the same BA transition.

No \texttt{ExternalSystem}, \texttt{TrustBoundary} or
\texttt{IntegrationRisk} metaclass is forced into BA merely because one
method projection wants it.

\subsection{16. Provider normalization
pressure}\label{16-provider-normalization-pressure}

Provider-state normalization remains a useful authority control.

A method projection that needs \texttt{PaymentAuthorizationResult} must
consume the accepted normalized BA meaning. It may add a method-local
interpretation such as an assurance state or decision branch, but it
must not bypass BA3 by mapping raw provider vocabulary privately into a
supposedly shared result.

Thus:

\begin{verbatim}
provider raw state
  -> governed mapping / BA3 derivation
  -> accepted BA PaymentAuthorizationResult
  -> method projection interpretation
\end{verbatim}

not:

\begin{verbatim}
provider raw state
  -> hidden method mapping
  -> shared project truth
\end{verbatim}

\subsection{17. Cross-projection comparison needs no new shared
identity}\label{17-cross-projection-comparison-needs-no-new-shared-identity}

T2 explicitly tries to force a common identity between method items.

The attempt fails under current evidence because BA trace already
provides the shared semantic anchor.

For same-baseline comparison:

\begin{verbatim}
item A trace set
  INTERSECT
item B trace set
\end{verbatim}

is sufficient to identify shared or overlapping BA basis.

For cross-baseline comparison, the BA3 continuity of the traced BA
elements supplies the semantic bridge before the projections are
rebuilt.

A persistent method-to-method equivalence relation would add consumer
taxonomy into the shared core without a demonstrated methodology-neutral
need.

Therefore:

\begin{verbatim}
SharedProjectionConcept identity                  NOT FORCED
method-category equivalence relation              NOT FORCED
new BA2 operator for projection comparison         NOT FORCED
\end{verbatim}

\subsection{18. Projection descriptor revision
changes}\label{18-projection-descriptor-revision-changes}

A projection descriptor revision is consumer evolution, not project
semantic change.

If the BA baseline is unchanged but the projection descriptor changes:

\begin{verbatim}
projection@descriptor-R1 remains historical
projection@descriptor-R2 is rebuilt from same BA
\end{verbatim}

Different local items do not imply that BA changed.

If a method wants continuity between its own projection items across
descriptor revisions, that remains method/tool-local history unless a
later concrete BA4 responsibility forces otherwise.

BA4-T2 still finds no need to duplicate BA3
\texttt{RETAIN\ \textbar{}\ REPLACE\ \textbar{}\ RETIRE} for projection
items.

\subsection{19. Reopen checks}\label{19-reopen-checks}

T2 finds no material counterexample to the closed BA0-BA3
responsibilities.

\begin{itemize}
\tightlist
\item
  BA1 still supplies all shared project-semantic identity required by
  both incompatible methods.
\item
  BA2 operators/roles preserve the method-neutral facts consumed by both
  projections.
\item
  BA3 supplies source provenance, review/freshness, continuity and
  change impact needed to qualify/rebuild projection input.
\end{itemize}

A method-specific taxonomy need is not evidence for BA reopen.

Reopen the smallest earlier responsibility only if a method cannot
obtain required \textbf{shared methodology-neutral project meaning} from
accepted BA without reconstructing it from governed prose or inventing
it privately.

No reviewed T2 case demonstrates that failure.

\subsection{20. Refined provisional
dispositions}\label{20-refined-provisional-dispositions}

\begin{verbatim}
T1 projection authority boundary                         PASS / RETAIN
T1 BA trace lower bound                                  PASS / RETAIN
T1 role-bound multi-input trace                          PASS / RETAIN
T1 generic selectionCoverageContract                     REFINED
Coverage mode distinction                                REQUIRED
  EXHAUSTIVE_FOR_DECLARED_SCOPE | SELECTIVE
Exhaustive eligible omission -> project absence           REJECTED
Exhaustive eligible omission -> projection defect         REQUIRED
Qualification policy                          REQUIRED
Diagnostic/stale review projection            PASS WITH QUALIFICATION
Non-current BA -> current shared projection truth         REJECTED
Descriptor + trace alone for non-trivial method rule      REJECTED
interpretationRuleRef                         REQUIRED FOR METHOD INTERPRETATION
Universal projection rule language                       REJECTED
Two incompatible method taxonomies                       PASS
Universal shared projection ontology                     REJECTED
Cross-projection taxonomy agreement                      NOT REQUIRED
Cross-projection comparison via BA trace                 PASS
New shared projection identity/relation                  NOT FORCED
Projection item BA-like lifecycle                        NOT FORCED
New BAE family                                           NOT FORCED
New BA2 operator                                         NOT FORCED
BA1 / BA2 / BA3 reopen                                   NOT TRIGGERED
BA4                                                      STARTED / NOT CLOSED
\end{verbatim}

\subsection{21. Falsification rules}\label{21-falsification-rules}

Revise this candidate if a concrete projection/corpus demonstrates that:

\begin{enumerate}
\def\labelenumi{\arabic{enumi}.}
\tightlist
\item
  non-trivial method-owned interpretation is reproducibly reviewable
  without identifying the applicable local rule;
\item
  selective and exhaustive-within-scope projections can be collapsed
  without losing omission/completeness semantics;
\item
  a stale/diagnostic projection can omit source qualification without
  being mistaken for current shared truth;
\item
  two incompatible method taxonomies require a shared projection
  ontology rather than BA trace;
\item
  cross-projection comparison requires a first-class shared identity not
  recoverable from BA identity/trace;
\item
  projection rebuild under M1-M4/order/provider requires a second
  project lifecycle model; or
\item
  a method consumer exposes missing shared methodology-neutral meaning
  that cannot be represented by closed BA1-BA3.
\end{enumerate}

No reviewed T2 evidence currently forces those revisions.

\subsection{22. Smallest unresolved set after
BA4-T2}\label{22-smallest-unresolved-set-after-ba4-t2}

The remaining BA4 work is now a distinct integrated closure review
rather than another taxonomy expansion.

The smallest unresolved set is:

\begin{enumerate}
\def\labelenumi{\arabic{enumi}.}
\tightlist
\item
  try to remove or merge \texttt{coverageMode}, qualification policy and
  \texttt{interpretationRuleRef} without reintroducing T2 failures;
\item
  verify that shared rendering, method interpretation and projection
  completeness remain orthogonal across both corpora;
\item
  regress T1/T2 projection rebuild and cross-projection comparison under
  M1-M4/order/WMS/provider;
\item
  verify once more that BA trace is sufficient as the common denominator
  without method-category equivalence; and
\item
  close BA4 for current thesis scope or reopen only the smallest failed
  BA4 responsibility.
\end{enumerate}

\subsection{23. Next falsification
target}\label{23-next-falsification-target}

Only after this package is reviewed, committed, pushed and remotely
verified:

\begin{quote}
\textbf{BA4-T3 - projection boundary, interpretation/coverage and
cross-projection closure review.}
\end{quote}

Do not start BA5 before BA4 is explicitly closed.

\end{document}
